% ============================================================
% CHAPITRE 21 — Illustration vectorielle et identité visuelle
% ============================================================

\chapter{Illustration vectorielle et identité visuelle}

\noindent
\textit{Dans le chapitre précédent, nous avons appris à créer des logos professionnels : comprendre la fonction d'un logo, analyser un brief, rechercher des références, esquisser, construire un logo géométrique ou typographique, évaluer la lisibilité, préparer les variantes. Le logo est l'élément central d'une identité visuelle. Mais une identité visuelle ne se limite pas à un logo : elle englobe l'ensemble des éléments graphiques qui permettent de reconnaître et de représenter une marque.}

\vspace{1em}

\noindent
L'\terme{identité visuelle} est bien plus qu'un logo : c'est un \emph{système cohérent} qui comprend le logo, les couleurs, les typographies, les éléments graphiques secondaires (icônes, pictogrammes, motifs, illustrations), les règles d'application, et les déclinaisons sur tous les supports (papeterie, packaging, web, signalétique, publicité). Une identité visuelle réussie est reconnaissable, cohérente, mémorable, intemporelle, adaptable. Une identité visuelle ratée est incohérente, générique, datée, difficile à décliner.

\vspace{1em}

\noindent
L'\terme{illustration vectorielle} joue un rôle central dans l'identité visuelle. Elle permet de créer des icônes, des pictogrammes, des illustrations, des motifs, qui enrichissent le système graphique et lui donnent de la personnalité. Une illustration vectorielle bien conçue s'intègre harmonieusement au logo, aux couleurs, aux typographies, et contribue à la reconnaissance de la marque. Une illustration mal conçue détonne, brouille le message, nuit à la cohérence.

\vspace{1em}

\noindent
Ce chapitre a pour objectif de te faire maîtriser l'illustration vectorielle et l'identité visuelle. Nous commencerons par les principes de l'illustration vectorielle. Nous apprendrons ensuite à créer des icônes, à concevoir des pictogrammes, à créer des illustrations simples, à harmoniser les formes et les couleurs, à développer un système graphique cohérent, à concevoir une palette de marque, à associer les couleurs et les typographies, à créer des éléments de papeterie, et à présenter une identité visuelle. À la fin de ce chapitre, tu sauras créer des identités visuelles complètes et professionnelles.

% ============================================================
\section{Les principes de l'illustration vectorielle}
% ============================================================

\subsection{Qu'est-ce que l'illustration vectorielle ?}

\begin{quote}
\textbf{L'illustration vectorielle est une illustration créée à partir de formes géométriques (tracés, courbes, points) décrites mathématiquement. Elle peut être agrandie ou réduite à l'infini sans perte de qualité.}
\end{quote}

L'illustration vectorielle se distingue de l'illustration matricielle (pixels) par sa nature mathématique. Elle est utilisée pour les logos, les icônes, les pictogrammes, les infographies, les illustrations techniques, les motifs.

\subsection{Les caractéristiques de l'illustration vectorielle}

\subsubsection{La précision}

L'illustration vectorielle est \terme{précise} : les formes sont nettes, les courbes sont fluides, les contours sont exacts.

\subsubsection{La scalabilité}

L'illustration vectorielle est \terme{scalable} : elle peut être agrandie ou réduite à l'infini sans perte de qualité.

\subsubsection{La légèreté}

L'illustration vectorielle est \terme{légère} : les fichiers sont souvent plus petits que les images matricielles.

\subsubsection{La modificabilité}

L'illustration vectorielle est \terme{modifiable} : chaque forme peut être éditée indépendamment.

\subsubsection{La cohérence}

L'illustration vectorielle est \terme{cohérente} : les formes partagent un style, une logique, une identité.

\subsection{Les styles d'illustration vectorielle}

\subsubsection{Le style flat}

Le \terme{style flat} utilise des formes pleines, sans contour, sans dégradé. Il est moderne, épuré, lisible.

\subsubsection{Le style ligne}

Le \terme{style ligne} (ou \emph{line art}) utilise des contours uniquement, sans remplissage. Il est élégant, léger, aéré.

\subsubsection{Le style plein}

Le \terme{style plein} utilise des formes remplies, avec ou sans contour. Il est fort, visible, impactant.

\subsubsection{Le style isométrique}

Le \terme{style isométrique} utilise une perspective isométrique (3D simulée). Il est technique, moderne, précis.

\subsubsection{Le style géométrique}

Le \terme{style géométrique} utilise des formes géométriques simples (cercles, carrés, triangles). Il est abstrait, minimaliste, intemporel.

\subsubsection{Le style organique}

Le \terme{style organique} utilise des formes courbes, fluides, naturelles. Il est chaleureux, humain, expressif.

\subsection{Les usages de l'illustration vectorielle}

\begin{itemize}
    \item \terme{Logos} : symbole, lettrage.
    \item \terme{Icônes} : interfaces, applications.
    \item \terme{Pictogrammes} : signalétique, infographies.
    \item \terme{Illustrations} : éditoriales, publicitaires.
    \item \terme{Motifs} : papeterie, packaging.
    \item \terme{Infographies} : graphiques, schémas.
    \item \terme{Typographie} : compositions, lettrage.
\end{itemize}

\subsection{Les bonnes pratiques}

\begin{itemize}
    \item \terme{Choisir} un style cohérent.
    \item \terme{Simplifier} les formes.
    \item \terme{Aligner} les éléments.
    \item \terme{Utiliser} une palette cohérente.
    \item \terme{Vérifier} la lisibilité à toutes les tailles.
    \item \terme{Organiser} les calques.
    \item \terme{Conserver} les fichiers sources.
\end{itemize}

\subsection*{Conclusion de la section}

L'\terme{illustration vectorielle} est une illustration créée à partir de formes géométriques. Elle est précise, scalable, légère, modifiable, cohérente. Il existe plusieurs styles : flat, ligne, plein, isométrique, géométrique, organique. Elle est utilisée pour les logos, les icônes, les pictogrammes, les illustrations, les motifs, les infographies, la typographie. Il est essentiel de choisir un style cohérent, de simplifier, d'aligner, de vérifier la lisibilité.

Maintenant que nous connaissons les principes, nous allons, dans la section suivante, apprendre à \terme{créer des icônes}.

% ============================================================
\section{Créer des icônes}
% ============================================================

\subsection{Qu'est-ce qu'une icône ?}

\begin{quote}
\textbf{Une icône est une représentation graphique simplifiée d'un objet, d'une action, d'un concept. Elle est utilisée dans les interfaces, la signalétique, les infographies.}
\end{quote}

Les icônes sont utilisées pour :

\begin{itemize}
    \item \terme{Les interfaces} : boutons, menus, notifications.
    \item \terme{Les applications} : fonctionnalités, actions.
    \item \terme{La signalétique} : panneaux, indications.
    \item \terme{Les infographies} : pictogrammes, symboles.
    \item \terme{Les sites web} : navigation, fonctionnalités.
\end{itemize}

\subsection{Les principes de création d'icônes}

\subsubsection{La simplicité}

Une icône est \terme{simple} : peu de détails, formes claires, lisible à petite taille.

\subsubsection{La reconnaissance}

Une icône est \terme{reconnaissable} : elle évoque immédiatement l'objet ou l'action.

\subsubsection{La cohérence}

Une icône est \terme{cohérente} avec les autres icônes : même style, même épaisseur, même palette.

\subsubsection{L'universalité}

Une icône est \terme{universelle} : comprise par tous, sans texte.

\subsection{Créer une icône}

\subsubsection{Étape 1 : Analyser}

\begin{itemize}
    \item \terme{Quel} est le sujet de l'icône ?
    \item \terme{Quelle} est sa fonction ?
    \item \terme{Quel} est le style souhaité (flat, ligne, plein) ?
\end{itemize}

\subsubsection{Étape 2 : Esquisser}

\begin{itemize}
    \item \terme{Dessiner} l'icône sur papier.
    \item \terme{Simplifier} au maximum.
    \item \terme{Tester} la reconnaissance.
\end{itemize}

\subsubsection{Étape 3 : Créer les formes de base}

\begin{itemize}
    \item \terme{Utiliser} des formes géométriques (cercle, carré, rectangle).
    \item \terme{Positionner} les formes.
    \item \terme{Aligner} les formes.
\end{itemize}

\subsubsection{Étape 4 : Dessiner les détails}

\begin{itemize}
    \item \terme{Dessiner} les courbes avec la Plume.
    \item \terme{Combiner} les formes (Pathfinder, Concepteur de forme).
    \item \terme{Simplifier} les tracés.
\end{itemize}

\subsubsection{Étape 5 : Appliquer les couleurs}

\begin{itemize}
    \item \terme{Appliquer} un remplissage (couleur unie).
    \item \terme{Appliquer} un contour (si nécessaire).
    \item \terme{Vérifier} la cohérence.
\end{itemize}

\subsubsection{Étape 6 : Vérifier}

\begin{itemize}
    \item \terme{Zoomer} pour vérifier à petite taille.
    \item \terme{Vérifier} la reconnaissance.
    \item \terme{Vérifier} la cohérence avec les autres icônes.
\end{itemize}

\subsection{Les styles d'icônes}

\subsubsection{Le style flat}

Les \terme{icônes flat} sont composées de formes pleines, sans contour. Elles sont modernes, lisibles.

\subsubsection{Le style ligne}

Les \terme{icônes ligne} sont composées de contours uniquement. Elles sont élégantes, légères.

\subsubsection{Le style plein}

Les \terme{icônes pleines} sont composées de formes remplies. Elles sont fortes, visibles.

\subsubsection{Le style mixte}

Les \terme{icônes mixtes} combinent formes pleines et contours. Elles sont riches, détaillées.

\subsection{Les bonnes pratiques}

\begin{itemize}
    \item \terme{Commencer} par esquisser.
    \item \terme{Simplifier} au maximum.
    \item \terme{Utiliser} des formes géométriques de base.
    \item \terme{Combiner} les formes avec le Pathfinder.
    \item \terme{Aligner} les éléments.
    \item \terme{Vérifier} la lisibilité à petite taille.
    \item \terme{Maintenir} la cohérence du style.
    \item \terme{Tester} sur différents fonds.
\end{itemize}

\subsection*{Conclusion de la section}

Les \terme{icônes} sont des représentations graphiques simplifiées, utilisées dans les interfaces, la signalétique, les infographies. Elles doivent être simples, reconnaissables, cohérentes, universelles. La création suit plusieurs étapes : analyser, esquisser, créer les formes de base, dessiner les détails, appliquer les couleurs, vérifier. Il est essentiel de commencer par esquisser, de simplifier, d'utiliser des formes géométriques, de vérifier la lisibilité.

Maintenant que nous savons créer des icônes, nous allons, dans la section suivante, apprendre à \terme{concevoir des pictogrammes}.

% ============================================================
\section{Concevoir des pictogrammes}
% ============================================================

\subsection{Qu'est-ce qu'un pictogramme ?}

\begin{quote}
\textbf{Un pictogramme est une représentation graphique simplifiée d'un objet, d'une action, d'un concept, destinée à être comprise universellement. Il est utilisé dans la signalétique, les infographies, les modes d'emploi.}
\end{quote}

Le pictogramme se distingue de l'icône par son usage : il est destiné à un public large, souvent dans un contexte de signalétique ou d'information. Il doit être compris sans texte, par tous.

\subsection{Les principes du pictogramme}

\subsubsection{La simplicité}

Un pictogramme est \terme{simple} : peu de détails, formes claires.

\subsubsection{L'universalité}

Un pictogramme est \terme{universel} : compris par tous, quelles que soient la langue, la culture, le niveau d'éducation.

\subsubsection{La reconnaissance}

Un pictogramme est \terme{reconnaissable} : il évoque immédiatement l'objet ou l'action.

\subsubsection{La cohérence}

Un pictogramme est \terme{cohérent} avec les autres pictogrammes : même style, même épaisseur, même logique.

\subsection{Les types de pictogrammes}

\subsubsection{Les pictogrammes de signalétique}

\begin{itemize}
    \item \terme{Direction} : flèches, indicateurs.
    \item \terme{Danger} : avertissements, interdictions.
    \item \terme{Obligation} : port du casque, ceinture.
    \item \terme{Information} : sortie, escalier, ascenseur.
\end{itemize}

\subsubsection{Les pictogrammes d'infographie}

\begin{itemize}
    \item \terme{Données} : statistiques, graphiques.
    \item \terme{Processus} : étapes, flux.
    \item \terme{Concepts} : idées, relations.
\end{itemize}

\subsubsection{Les pictogrammes d'interface}

\begin{itemize}
    \item \terme{Actions} : boutons, commandes.
    \item \terme{États} : succès, erreur, attention.
    \item \terme{Navigation} : menus, liens.
\end{itemize}

\subsection{Concevoir un pictogramme}

\subsubsection{Étape 1 : Analyser}

\begin{itemize}
    \item \terme{Quel} est le sujet du pictogramme ?
    \item \terme{Quelle} est sa fonction ?
    \item \terme{Quel} est le public ?
\end{itemize}

\subsubsection{Étape 2 : Esquisser}

\begin{itemize}
    \item \terme{Dessiner} le pictogramme sur papier.
    \item \terme{Simplifier} au maximum.
    \item \terme{Tester} la reconnaissance.
\end{itemize}

\subsubsection{Étape 3 : Créer}

\begin{itemize}
    \item \terme{Utiliser} des formes géométriques.
    \item \terme{Dessiner} avec la Plume.
    \item \terme{Combiner} les formes.
    \item \terme{Simplifier} les tracés.
\end{itemize}

\subsubsection{Étape 4 : Vérifier}

\begin{itemize}
    \item \terme{Vérifier} la reconnaissance.
    \item \terme{Vérifier} l'universalité.
    \item \terme{Vérifier} la cohérence.
\end{itemize}

\subsection{Les bonnes pratiques}

\begin{itemize}
    \item \terme{Commencer} par esquisser.
    \item \terme{Simplifier} au maximum.
    \item \terme{Utiliser} des formes géométriques.
    \item \terme{Vérifier} la reconnaissance.
    \item \terme{Vérifier} l'universalité.
    \item \terme{Maintenir} la cohérence.
    \item \terme{Tester} sur différents supports.
\end{itemize}

\subsection*{Conclusion de la section}

Les \terme{pictogrammes} sont des représentations graphiques simplifiées, destinées à être comprises universellement. Ils sont utilisés dans la signalétique, les infographies, les interfaces. Ils doivent être simples, universels, reconnaissables, cohérents. La conception suit plusieurs étapes : analyser, esquisser, créer, vérifier. Il est essentiel de commencer par esquisser, de simplifier, de vérifier la reconnaissance.

Maintenant que nous savons concevoir des pictogrammes, nous allons, dans la section suivante, apprendre à \terme{créer des illustrations simples}.

% ============================================================
\section{Créer des illustrations simples}
% ============================================================

\subsection{Qu'est-ce qu'une illustration simple ?}

\begin{quote}
\textbf{Une illustration simple est une image vectorielle qui représente un objet, un concept, une scène, avec un niveau de détail modéré. Elle est utilisée dans l'édition, la publicité, le web, les identités visuelles.}
\end{quote}

Les illustrations simples sont utilisées pour :

\begin{itemize}
    \item \terme{Les identités visuelles} : illustrations de marque.
    \item \terme{Les sites web} : illustrations d'interface.
    \item \terme{Les publicités} : illustrations publicitaires.
    \item \terme{Les éditions} : illustrations éditoriales.
    \item \terme{Les infographies} : illustrations explicatives.
\end{itemize}

\subsection{Les principes de l'illustration simple}

\subsubsection{La simplicité}

Une illustration simple est \terme{simple} : peu de formes, peu de détails, peu de couleurs.

\subsubsection{La cohérence}

Une illustration simple est \terme{cohérente} avec les autres éléments de l'identité visuelle : même style, même palette, même logique.

\subsubsection{L'expressivité}

Une illustration simple est \terme{expressive} : elle transmet une ambiance, une émotion, une idée.

\subsubsection{La lisibilité}

Une illustration simple est \terme{lisible} : elle est comprise immédiatement.

\subsection{Créer une illustration simple}

\subsubsection{Étape 1 : Concevoir}

\begin{itemize}
    \item \terme{Définir} le sujet de l'illustration.
    \item \terme{Définir} l'ambiance.
    \item \terme{Esquisser} plusieurs pistes.
\end{itemize}

\subsubsection{Étape 2 : Construire les formes de base}

\begin{itemize}
    \item \terme{Utiliser} des formes géométriques.
    \item \terme{Positionner} les formes.
    \item \terme{Aligner} les formes.
\end{itemize}

\subsubsection{Étape 3 : Dessiner les détails}

\begin{itemize}
    \item \terme{Dessiner} avec la Plume.
    \item \terme{Combiner} les formes.
    \item \terme{Simplifier} les tracés.
\end{itemize}

\subsubsection{Étape 4 : Appliquer les couleurs}

\begin{itemize}
    \item \terme{Choisir} une palette cohérente.
    \item \terme{Appliquer} les couleurs.
    \item \terme{Vérifier} le contraste.
\end{itemize}

\subsubsection{Étape 5 : Vérifier}

\begin{itemize}
    \item \terme{Vérifier} la lisibilité.
    \item \terme{Vérifier} la cohérence.
    \item \terme{Vérifier} l'expressivité.
\end{itemize}

\subsection{Les styles d'illustration simple}

\subsubsection{Le style flat}

Le \terme{style flat} utilise des formes pleines, sans contour, sans dégradé. Il est moderne, épuré.

\subsubsection{Le style ligne}

Le \terme{style ligne} utilise des contours uniquement. Il est élégant, léger.

\subsubsection{Le style géométrique}

Le \terme{style géométrique} utilise des formes géométriques simples. Il est abstrait, minimaliste.

\subsubsection{Le style organique}

Le \terme{style organique} utilise des formes courbes, fluides. Il est chaleureux, humain.

\subsection{Les bonnes pratiques}

\begin{itemize}
    \item \terme{Commencer} par esquisser.
    \item \terme{Utiliser} des formes géométriques de base.
    \item \terme{Simplifier} au maximum.
    \item \terme{Utiliser} une palette cohérente.
    \item \terme{Vérifier} la lisibilité.
    \item \terme{Maintenir} la cohérence avec l'identité visuelle.
    \item \terme{Tester} sur différents supports.
\end{itemize}

\subsection*{Conclusion de la section}

Les \terme{illustrations simples} sont des images vectorielles qui représentent un objet, un concept, une scène. Elles sont utilisées dans les identités visuelles, les sites web, les publicités, les éditions, les infographies. Elles doivent être simples, cohérentes, expressives, lisibles. La création suit plusieurs étapes : concevoir, construire les formes de base, dessiner les détails, appliquer les couleurs, vérifier. Il est essentiel de commencer par esquisser, de simplifier, de vérifier la lisibilité.

Maintenant que nous savons créer des illustrations simples, nous allons, dans la section suivante, apprendre à \terme{harmoniser les formes et les couleurs}.

% ============================================================
\section{Harmoniser les formes et les couleurs}
% ============================================================

\subsection{Pourquoi harmoniser ?}

L'\terme{harmonisation} des formes et des couleurs est essentielle pour :

\begin{itemize}
    \item \terme{Créer} un ensemble cohérent.
    \item \terme{Assurer} la reconnaissance.
    \item \terme{Faciliter} la lecture.
    \item \terme{Donner} une image professionnelle.
    \item \terme{Respecter} l'identité visuelle.
\end{itemize}

Une identité visuelle harmonieuse est une identité visuelle où tous les éléments (logo, icônes, pictogrammes, illustrations, couleurs, typographies) s'accordent pour former un tout.

\subsection{L'harmonisation des formes}

\subsubsection{La cohérence des formes}

Les formes doivent être \terme{cohérentes} :

\begin{itemize}
    \item \terme{Même} style (flat, ligne, plein).
    \item \terme{Même} épaisseur de contour.
    \item \terme{Même} niveau de détail.
    \item \terme{Même} logique géométrique.
\end{itemize}

\subsubsection{La simplicité des formes}

Les formes doivent être \terme{simples} :

\begin{itemize}
    \item \terme{Éviter} les détails inutiles.
    \item \terme{Éviter} les formes trop complexes.
    \item \terme{Éviter} les effets superflus.
\end{itemize}

\subsubsection{La répétition des formes}

Les formes peuvent être \terme{répétées} :

\begin{itemize}
    \item \terme{Répéter} un motif.
    \item \terme{Répéter} une forme de base.
    \item \terme{Répéter} une logique.
\end{itemize}

\subsection{L'harmonisation des couleurs}

\subsubsection{La palette de couleurs}

La \terme{palette} doit être cohérente :

\begin{itemize}
    \item \terme{Limiter} le nombre de couleurs (3 à 5).
    \item \terme{Choisir} une harmonie (monochromatique, analogue, complémentaire).
    \item \terme{Définir} une couleur principale, des couleurs secondaires, des accents.
\end{itemize}

\subsubsection{La cohérence des couleurs}

Les couleurs doivent être \terme{cohérentes} :

\begin{itemize}
    \item \terme{Utiliser} les mêmes couleurs dans tous les éléments.
    \item \terme{Respecter} les valeurs (clair/foncé).
    \item \terme{Respecter} les saturations (vif/terne).
\end{itemize}

\subsubsection{Le contraste}

Le \terme{contraste} doit être suffisant :

\begin{itemize}
    \item \terme{Contraste} entre texte et fond.
    \item \terme{Contraste} entre éléments.
    \item \terme{Contraste} entre couleurs.
\end{itemize}

\subsection{L'harmonisation entre formes et couleurs}

\subsubsection{Le principe}

Les formes et les couleurs doivent \terme{s'harmoniser} :

\begin{itemize}
    \item \terme{Les} formes simples appellent des couleurs simples.
    \item \terme{Les} formes complexes appellent des couleurs sobres.
    \item \terme{Les} formes organiques appellent des couleurs chaudes.
    \item \terme{Les} formes géométriques appellent des couleurs froides.
\end{itemize}

\subsubsection{Les bonnes pratiques}

\begin{itemize}
    \item \terme{Choisir} des couleurs cohérentes.
    \item \terme{Choisir} des formes cohérentes.
    \item \terme{Vérifier} l'harmonie globale.
    \item \terme{Tester} sur différents supports.
\end{itemize}

\subsection{Les bonnes pratiques}

\begin{itemize}
    \item \terme{Définir} une palette cohérente.
    \item \terme{Définir} un style de formes.
    \item \terme{Répéter} les formes.
    \item \terme{Répéter} les couleurs.
    \item \terme{Vérifier} le contraste.
    \item \terme{Vérifier} l'harmonie.
    \item \terme{Tester} sur différents supports.
\end{itemize}

\subsection*{Conclusion de la section}

L'\terme{harmonisation} des formes et des couleurs est essentielle pour créer un ensemble cohérent. Les formes doivent être cohérentes, simples, répétées. Les couleurs doivent être cohérentes, contrastées, harmonieuses. Les formes et les couleurs doivent s'harmoniser. Il est essentiel de définir une palette, un style de formes, de vérifier l'harmonie, de tester sur différents supports.

Maintenant que nous savons harmoniser, nous allons, dans la section suivante, apprendre à \terme{développer un système graphique cohérent}.

% ============================================================
\section{Développer un système graphique cohérent}
% ============================================================

\subsection{Qu'est-ce qu'un système graphique ?}

\begin{quote}
\textbf{Un système graphique est un ensemble cohérent d'éléments visuels (logo, couleurs, typographies, icônes, illustrations, motifs, règles) qui permettent de représenter une marque sur tous les supports.}
\end{quote}

Le système graphique est le \terme{squelette} de l'identité visuelle. Il définit :

\begin{itemize}
    \item \terme{Les} éléments de base (logo, couleurs, typographies).
    \item \terme{Les} éléments secondaires (icônes, pictogrammes, illustrations, motifs).
    \item \terme{Les} règles d'application.
    \item \terme{Les} déclinaisons.
\end{itemize}

\subsection{Les composants d'un système graphique}

\subsubsection{Le logo}

Le \terme{logo} est l'élément central. Il définit l'identité de la marque.

\subsubsection{Les couleurs}

Les \terme{couleurs} définissent l'ambiance. Elles sont organisées en palette (principale, secondaires, accents).

\subsubsection{Les typographies}

Les \terme{typographies} définissent la personnalité. Elles sont organisées en familles (titres, texte courant, accents).

\subsubsection{Les icônes et pictogrammes}

Les \terme{icônes} et \terme{pictogrammes} enrichissent le système. Ils sont cohérents avec le style du logo.

\subsubsection{Les illustrations}

Les \terme{illustrations} donnent de la personnalité. Elles sont cohérentes avec le style général.

\subsubsection{Les motifs}

Les \terme{motifs} ajoutent de la texture. Ils sont utilisés pour les fonds, les packagings, les papeteries.

\subsubsection{Les règles d'application}

Les \terme{règles} définissent comment utiliser les éléments. Elles sont réunies dans une charte graphique.

\subsection{Développer un système graphique}

\subsubsection{Étape 1 : Définir l'identité}

\begin{itemize}
    \item \terme{Définir} les valeurs de la marque.
    \item \terme{Définir} la personnalité.
    \item \terme{Définir} le positionnement.
\end{itemize}

\subsubsection{Étape 2 : Créer les éléments de base}

\begin{itemize}
    \item \terme{Créer} le logo.
    \item \terme{Définir} la palette.
    \item \terme{Choisir} les typographies.
\end{itemize}

\subsubsection{Étape 3 : Créer les éléments secondaires}

\begin{itemize}
    \item \terme{Créer} les icônes.
    \item \terme{Créer} les pictogrammes.
    \item \terme{Créer} les illustrations.
    \item \terme{Créer} les motifs.
\end{itemize}

\subsubsection{Étape 4 : Définir les règles}

\begin{itemize}
    \item \terme{Définir} les usages du logo.
    \item \terme{Définir} les usages des couleurs.
    \item \terme{Définir} les usages des typographies.
    \item \terme{Définir} les usages des éléments secondaires.
\end{itemize}

\subsubsection{Étape 5 : Décliner}

\begin{itemize}
    \item \terme{Décliner} sur la papeterie.
    \item \terme{Décliner} sur le web.
    \item \terme{Décliner} sur le packaging.
    \item \terme{Décliner} sur la signalétique.
    \item \terme{Décliner} sur la publicité.
\end{itemize}

\subsubsection{Étape 6 : Documenter}

\begin{itemize}
    \item \terme{Créer} une charte graphique.
    \item \terme{Documenter} les règles.
    \item \terme{Présenter} les exemples.
\end{itemize}

\subsection{Les bonnes pratiques}

\begin{itemize}
    \item \terme{Définir} une identité claire.
    \item \terme{Créer} des éléments cohérents.
    \item \terme{Définir} des règles précises.
    \item \terme{Décliner} sur tous les supports.
    \item \terme{Documenter} dans une charte.
    \item \terme{Tester} sur différents supports.
    \item \terme{Faire} évoluer le système.
\end{itemize}

\subsection*{Conclusion de la section}

Un \terme{système graphique} est un ensemble cohérent d'éléments visuels. Il comprend le logo, les couleurs, les typographies, les icônes, les pictogrammes, les illustrations, les motifs, les règles. Le développement suit plusieurs étapes : définir l'identité, créer les éléments de base, créer les éléments secondaires, définir les règles, décliner, documenter. Il est essentiel de définir une identité claire, de créer des éléments cohérents, de décliner, de documenter.

Maintenant que nous savons développer un système graphique, nous allons, dans la section suivante, apprendre à \terme{concevoir une palette de marque}.

% ============================================================
\section{Concevoir une palette de marque}
% ============================================================

\subsection{Qu'est-ce qu'une palette de marque ?}

\begin{quote}
\textbf{Une palette de marque est un ensemble limité de couleurs définies pour une marque, organisées en couleurs principales, secondaires et accents. Elle définit l'ambiance et la personnalité de la marque.}
\end{quote}

La palette de marque est essentielle pour :

\begin{itemize}
    \item \terme{Créer} une ambiance cohérente.
    \item \terme{Assurer} la reconnaissance.
    \item \terme{Transmettre} des valeurs.
    \item \terme{Différencier} la marque.
    \item \terme{Faciliter} les déclinaisons.
\end{itemize}

\subsection{Les composants d'une palette}

\subsubsection{Les couleurs principales}

Les \terme{couleurs principales} sont les couleurs dominantes. Elles définissent l'ambiance. Elles sont généralement au nombre de 1 à 3.

\subsubsection{Les couleurs secondaires}

Les \terme{couleurs secondaires} complètent les couleurs principales. Elles sont utilisées pour les éléments importants, mais pas dominants.

\subsubsection{Les couleurs d'accent}

Les \terme{couleurs d'accent} sont utilisées pour attirer l'attention sur des éléments spécifiques (boutons, liens, titres).

\subsubsection{Les couleurs neutres}

Les \terme{couleurs neutres} (blanc, noir, gris) servent de base, de fond, de texte.

\subsection{Concevoir une palette}

\subsubsection{Étape 1 : Analyser}

\begin{itemize}
    \item \terme{Quelles} sont les valeurs de la marque ?
    \item \terme{Quelle} est la personnalité ?
    \item \terme{Quel} est le positionnement ?
    \item \terme{Quel} est le public ?
\end{itemize}

\subsubsection{Étape 2 : Choisir une couleur principale}

\begin{itemize}
    \item \terme{Choisir} une couleur qui correspond à l'identité.
    \item \terme{Vérifier} les significations culturelles.
    \item \terme{Vérifier} la différenciation par rapport aux concurrents.
\end{itemize}

\subsubsection{Étape 3 : Choisir les couleurs secondaires}

\begin{itemize}
    \item \terme{Choisir} des couleurs qui complètent la couleur principale.
    \item \terme{Utiliser} une harmonie (analogue, complémentaire, triadique).
    \item \terme{Vérifier} la cohérence.
\end{itemize}

\subsubsection{Étape 4 : Choisir les couleurs d'accent}

\begin{itemize}
    \item \terme{Choisir} une couleur vive, contrastée.
    \item \terme{Utiliser} avec parcimonie.
    \item \terme{Vérifier} la lisibilité.
\end{itemize}

\subsubsection{Étape 5 : Choisir les couleurs neutres}

\begin{itemize}
    \item \terme{Choisir} du blanc, du noir, du gris.
    \item \terme{Définir} les usages (fond, texte, bordures).
\end{itemize}

\subsubsection{Étape 6 : Tester}

\begin{itemize}
    \item \terme{Tester} la palette sur différents supports.
    \item \terme{Vérifier} la cohérence.
    \item \terme{Vérifier} la lisibilité.
    \item \terme{Vérifier} le contraste.
\end{itemize}

\subsection{Les bonnes pratiques}

\begin{itemize}
    \item \terme{Limiter} le nombre de couleurs (3 à 5).
    \item \terme{Définir} une couleur principale.
    \item \terme{Utiliser} une harmonie.
    \item \terme{Définir} les usages.
    \item \terme{Tester} sur différents supports.
    \item \terme{Vérifier} les significations.
    \item \terme{Vérifier} la différenciation.
\end{itemize}

\subsection*{Conclusion de la section}

Une \terme{palette de marque} est un ensemble limité de couleurs définies pour une marque. Elle comprend des couleurs principales, secondaires, accents, neutres. La conception suit plusieurs étapes : analyser, choisir la couleur principale, choisir les couleurs secondaires, choisir les accents, choisir les neutres, tester. Il est essentiel de limiter le nombre de couleurs, de définir une couleur principale, de tester sur différents supports.

Maintenant que nous savons concevoir une palette, nous allons, dans la section suivante, apprendre à \terme{associer les couleurs et les typographies}.

% ============================================================
\section{Associer les couleurs et les typographies}
% ============================================================

\subsection{Pourquoi associer couleurs et typographies ?}

L'\terme{association} des couleurs et des typographies est essentielle pour :

\begin{itemize}
    \item \terme{Créer} une ambiance cohérente.
    \item \terme{Assurer} la lisibilité.
    \item \terme{Transmettre} une personnalité.
    \item \terme{Différencier} la marque.
    \item \terme{Respecter} l'identité visuelle.
\end{itemize}

Une association réussie renforce le message ; une association ratée le brouille.

\subsection{Les principes de l'association}

\subsubsection{La cohérence}

Les couleurs et les typographies doivent être \terme{cohérentes} :

\begin{itemize}
    \item \terme{Même} ambiance (élégante, moderne, créative).
    \item \terme{Même} niveau de contraste.
    \item \terme{Même} personnalité.
\end{itemize}

\subsubsection{La lisibilité}

L'association doit être \terme{lisible} :

\begin{itemize}
    \item \terme{Contraste} suffisant entre texte et fond.
    \item \terme{Contraste} suffisant entre titres et texte.
    \item \terme{Cohérence} des couleurs.
\end{itemize}

\subsubsection{L'harmonie}

L'association doit être \terme{harmonieuse} :

\begin{itemize}
    \item \terme{Couleurs} qui s'accordent.
    \item \terme{Typographies} qui s'accordent.
    \item \terme{Ensemble} équilibré.
\end{itemize}

\subsection{Les bonnes associations}

\subsubsection{Couleurs sombres et typographies claires}

Les \terme{couleurs sombres} (noir, bleu foncé, gris anthracite) s'associent bien avec des \terme{typographies claires} (blanc, crème). Cette association est élégante, lisible.

\subsubsection{Couleurs claires et typographies sombres}

Les \terme{couleurs claires} (blanc, beige, pastel) s'associent bien avec des \terme{typographies sombres} (noir, gris foncé). Cette association est moderne, épurée.

\subsubsection{Couleurs vives et typographies sobres}

Les \terme{couleurs vives} (rouge, orange, jaune) s'associent bien avec des \terme{typographies sobres} (blanc, noir). Cette association est dynamique, impactante.

\subsubsection{Couleurs neutres et typographies contrastées}

Les \terme{couleurs neutres} (gris, beige) s'associent bien avec des \terme{typographies contrastées} (noir, blanc). Cette association est sobre, professionnelle.

\subsection{Les erreurs à éviter}

\subsubsection{Couleurs et typographies trop similaires}

Des couleurs et des typographies trop similaires créent un manque de contraste, une perte de lisibilité.

\subsubsection{Couleurs et typographies trop différentes}

Des couleurs et des typographies trop différentes créent un choc visuel, une incohérence.

\subsubsection{Couleurs et typographies inadaptées au message}

Des couleurs et des typographies inadaptées au message brouillent la communication.

\subsubsection{Manque de contraste}

Un manque de contraste entre texte et fond rend le texte illisible.

\subsection{Les bonnes pratiques}

\begin{itemize}
    \item \terme{Définir} une palette cohérente.
    \item \terme{Définir} des typographies cohérentes.
    \item \terme{Vérifier} le contraste.
    \item \terme{Vérifier} la lisibilité.
    \item \terme{Vérifier} l'harmonie.
    \item \terme{Tester} sur différents supports.
    \item \terme{Documenter} dans la charte.
\end{itemize}

\subsection*{Conclusion de la section}

L'\terme{association} des couleurs et des typographies est essentielle pour créer une ambiance cohérente. Elle repose sur la cohérence, la lisibilité, l'harmonie. Les bonnes associations sont : couleurs sombres et typographies claires, couleurs claires et typographies sombres, couleurs vives et typographies sobres, couleurs neutres et typographies contrastées. Il est essentiel de définir une palette et des typographies cohérentes, de vérifier le contraste, de tester sur différents supports.

Maintenant que nous savons associer couleurs et typographies, nous allons, dans la section suivante, apprendre à \terme{créer des éléments de papeterie}.

% ============================================================
\section{Créer des éléments de papeterie}
% ============================================================

\subsection{Qu'est-ce que la papeterie ?}

\begin{quote}
\textbf{La papeterie est l'ensemble des supports imprimés d'une entreprise ou d'une marque : carte de visite, papier à en-tête, enveloppe, pochette, chemise, bloc-notes.}
\end{quote}

La papeterie est essentielle pour :

\begin{itemize}
    \item \terme{Communiquer} l'identité de la marque.
    \item \terme{Créer} une impression professionnelle.
    \item \terme{Faciliter} les échanges commerciaux.
    \item \terme{Marquer} les esprits.
    \item \terme{Assurer} la cohérence.
\end{itemize}

\subsection{Les éléments de papeterie}

\subsubsection{La carte de visite}

La \terme{carte de visite} est le support le plus personnel. Elle contient :

\begin{itemize}
    \item \terme{Le} logo.
    \item \terme{Le} nom.
    \item \terme{La} fonction.
    \item \terme{Les} coordonnées (téléphone, email, adresse, site web).
    \item \terme{Les} réseaux sociaux.
\end{itemize}

Format standard : 85 × 55 mm.

\subsubsection{Le papier à en-tête}

Le \terme{papier à en-tête} est utilisé pour les courriers officiels. Il contient :

\begin{itemize}
    \item \terme{Le} logo.
    \item \terme{Les} coordonnées.
    \item \terme{Les} mentions légales.
    \item \terme{Un} espace pour le contenu.
\end{itemize}

Format standard : A4 (210 × 297 mm).

\subsubsection{L'enveloppe}

L'\terme{enveloppe} est utilisée pour les courriers. Elle contient :

\begin{itemize}
    \item \terme{Le} logo.
    \item \terme{L'}adresse de l'expéditeur.
    \item \terme{Un} espace pour l'adresse du destinataire.
\end{itemize}

Formats standards : DL (110 × 220 mm), C4 (229 × 324 mm), C5 (162 × 229 mm).

\subsubsection{La pochette}

La \terme{pochette} (ou \emph{chemise}) est utilisée pour les documents. Elle contient :

\begin{itemize}
    \item \terme{Le} logo.
    \item \terme{Les} coordonnées.
    \item \terme{Un} espace pour le titre du document.
\end{itemize}

Format standard : A4.

\subsubsection{Le bloc-notes}

Le \terme{bloc-notes} est un objet promotionnel. Il contient :

\begin{itemize}
    \item \terme{Le} logo.
    \item \terme{Les} coordonnées.
    \item \terme{Un} espace pour écrire.
\end{itemize}

Format standard : A5 ou A6.

\subsection{Créer des éléments de papeterie}

\subsubsection{Étape 1 : Définir les éléments}

\begin{itemize}
    \item \terme{Le} logo (version adaptée).
    \item \terme{Les} couleurs de la marque.
    \item \terme{Les} typographies.
    \item \terme{Les} coordonnées.
\end{itemize}

\subsubsection{Étape 2 : Composer}

\begin{itemize}
    \item \terme{Placer} le logo.
    \item \terme{Placer} les coordonnées.
    \item \terme{Respecter} les marges.
    \item \terme{Respecter} la hiérarchie.
    \item \terme{Respecter} la grille.
\end{itemize}

\subsubsection{Étape 3 : Vérifier}

\begin{itemize}
    \item \terme{Vérifier} la lisibilité.
    \item \terme{Vérifier} la cohérence.
    \item \terme{Vérifier} la qualité.
\end{itemize}

\subsubsection{Étape 4 : Préparer l'impression}

\begin{itemize}
    \item \terme{Convertir} en CMJN.
    \item \terme{Prévoir} les fonds perdus.
    \item \terme{Ajouter} les traits de coupe.
    \item \terme{Incorporer} les polices.
    \item \terme{Exporter} en PDF/X.
\end{itemize}

\subsection{Les bonnes pratiques}

\begin{itemize}
    \item \terme{Utiliser} les éléments de l'identité.
    \item \terme{Respecter} les marges.
    \item \terme{Respecter} la grille.
    \item \terme{Vérifier} la lisibilité.
    \item \terme{Vérifier} la cohérence.
    \item \terme{Préparer} l'impression.
    \item \terme{Tester} sur papier.
\end{itemize}

\subsection*{Conclusion de la section}

La \terme{papeterie} est l'ensemble des supports imprimés d'une marque : carte de visite, papier à en-tête, enveloppe, pochette, bloc-notes. Elle est essentielle pour communiquer l'identité, créer une impression professionnelle, faciliter les échanges. La création suit plusieurs étapes : définir les éléments, composer, vérifier, préparer l'impression. Il est essentiel d'utiliser les éléments de l'identité, de respecter les marges, de vérifier la lisibilité, de préparer l'impression.

Maintenant que nous savons créer de la papeterie, nous allons, dans la dernière section de ce chapitre, apprendre à \terme{présenter une identité visuelle}.

% ============================================================
\section{Présenter une identité visuelle}
% ============================================================

\subsection{Pourquoi présenter une identité visuelle ?}

La \terme{présentation} d'une identité visuelle est essentielle pour :

\begin{itemize}
    \item \terme{Communiquer} la démarche.
    \item \terme{Expliquer} les choix.
    \item \terme{Convaincre} le client.
    \item \terme{Démontrer} la cohérence.
    \item \terme{Faciliter} la validation.
    \item \terme{Documenter} le projet.
\end{itemize}

Une identité visuelle bien présentée est une identité visuelle comprise, acceptée, valorisée.

\subsection{Les éléments de la présentation}

\subsubsection{La page de couverture}

La \terme{couverture} présente le projet :

\begin{itemize}
    \item \terme{Le} nom de la marque.
    \item \terme{Le} titre du projet.
    \item \terme{Le} logo.
    \item \terme{La} date.
\end{itemize}

\subsubsection{Le brief}

Le \terme{brief} rappelle le contexte :

\begin{itemize}
    \item \terme{Les} objectifs.
    \item \terme{Le} public.
    \item \terme{Les} contraintes.
\end{itemize}

\subsubsection{La démarche}

La \terme{démarche} explique les choix :

\begin{itemize}
    \item \terme{La} recherche.
    \item \terme{Les} inspirations.
    \item \terme{Les} esquisses.
    \item \terme{Les} concepts.
\end{itemize}

\subsubsection{Le logo}

Le \terme{logo} est présenté avec :

\begin{itemize}
    \item \terme{Les} variantes.
    \item \terme{Les} versions monochromes.
    \item \terme{Les} zones de protection.
    \item \terme{Les} usages.
\end{itemize}

\subsubsection{Les couleurs}

Les \terme{couleurs} sont présentées avec :

\begin{itemize}
    \item \terme{La} palette.
    \item \terme{Les} codes (RVB, CMJN, Pantone).
    \item \terme{Les} usages.
\end{itemize}

\subsubsection{Les typographies}

Les \terme{typographies} sont présentées avec :

\begin{itemize}
    \item \terme{Les} familles.
    \item \terme{Les} usages (titres, texte courant).
    \item \terme{Les} exemples.
\end{itemize}

\subsubsection{Les applications}

Les \terme{applications} présentent le système :

\begin{itemize}
    \item \terme{La} papeterie.
    \item \terme{Le} web.
    \item \terme{Le} packaging.
    \item \terme{La} signalétique.
    \item \terme{La} publicité.
\end{itemize}

\subsection{Présenter l'identité visuelle}

\subsubsection{Étape 1 : Structurer}

\begin{itemize}
    \item \terme{Organiser} la présentation.
    \item \terme{Définir} les sections.
    \item \terme{Créer} une cohérence.
\end{itemize}

\subsubsection{Étape 2 : Rédiger}

\begin{itemize}
    \item \terme{Expliquer} les choix.
    \item \terme{Justifier} les décisions.
    \item \terme{Démontrer} la cohérence.
\end{itemize}

\subsubsection{Étape 3 : Illustrer}

\begin{itemize}
    \item \terme{Présenter} les visuels.
    \item \terme{Montrer} les applications.
    \item \terme{Utiliser} des mockups.
\end{itemize}

\subsubsection{Étape 4 : Présenter}

\begin{itemize}
    \item \terme{Présenter} oralement.
    \item \terme{Expliquer} la démarche.
    \item \terme{Répondre} aux questions.
\end{itemize}

\subsection{Les bonnes pratiques}

\begin{itemize}
    \item \terme{Structurer} la présentation.
    \item \terme{Expliquer} les choix.
    \item \terme{Justifier} les décisions.
    \item \terme{Illustrer} avec des visuels.
    \item \terme{Présenter} les applications.
    \item \terme{Utiliser} des mockups.
    \item \terme{Documenter} la charte.
    \item \terme{Présenter} oralement.
\end{itemize}

\subsection*{Conclusion de la section}

La \terme{présentation} d'une identité visuelle est essentielle pour communiquer la démarche, expliquer les choix, convaincre le client, démontrer la cohérence. Elle comprend la couverture, le brief, la démarche, le logo, les couleurs, les typographies, les applications. La présentation suit plusieurs étapes : structurer, rédiger, illustrer, présenter. Il est essentiel de structurer, d'expliquer, d'illustrer, de présenter oralement.

% ============================================================
\section*{Conclusion du chapitre}
\addcontentsline{toc}{section}{Conclusion du chapitre}
% ============================================================

Nous voici au terme du chapitre 21. Nous avons appris à maîtriser l'illustration vectorielle et l'identité visuelle :

\begin{itemize}
    \item \terme{Principes de l'illustration vectorielle} : précision, scalabilité, légèreté, modificabilité, cohérence ; styles (flat, ligne, plein, isométrique, géométrique, organique).
    \item \terme{Création d'icônes} : simplicité, reconnaissance, cohérence, universalité ; méthode.
    \item \terme{Conception de pictogrammes} : simplicité, universalité, reconnaissance, cohérence ; types.
    \item \terme{Création d'illustrations simples} : simplicité, cohérence, expressivité, lisibilité ; méthode.
    \item \terme{Harmonisation des formes et des couleurs} : cohérence, simplicité, répétition ; palette, contraste.
    \item \terme{Développement d'un système graphique} : composants, méthode.
    \item \terme{Conception d'une palette de marque} : couleurs principales, secondaires, accents, neutres ; méthode.
    \item \terme{Association couleurs/typographies} : cohérence, lisibilité, harmonie ; bonnes associations, erreurs.
    \item \terme{Création d'éléments de papeterie} : carte de visite, papier à en-tête, enveloppe, pochette, bloc-notes ; méthode.
    \item \terme{Présentation d'une identité visuelle} : couverture, brief, démarche, logo, couleurs, typographies, applications ; méthode.
\end{itemize}

L'identité visuelle est bien plus qu'un logo : c'est un système cohérent qui comprend le logo, les couleurs, les typographies, les éléments graphiques secondaires, les règles d'application, et les déclinaisons sur tous les supports. L'illustration vectorielle joue un rôle central dans ce système : elle permet de créer des icônes, des pictogrammes, des illustrations, des motifs, qui enrichissent l'identité et lui donnent de la personnalité.

Une identité visuelle réussie est reconnaissable, cohérente, mémorable, intemporelle, adaptable. Une identité visuelle ratée est incohérente, générique, datée, difficile à décliner. La différence entre les deux repose sur la méthode, la rigueur, la cohérence, la précision.

Ce chapitre t'a donné les bases de l'illustration vectorielle et de l'identité visuelle. Mais la maîtrise vient avec la pratique. N'hésite pas à t'entraîner sur des projets variés : marques fictives, refontes d'identités, exercices créatifs. Chaque projet est un nouveau défi, et chaque défi est une occasion de progresser.

Ce chapitre clôt la troisième partie de ce cours, consacrée à la \terme{maîtrise d'Adobe Illustrator}. Tu disposes désormais des compétences nécessaires pour créer des logos, des icônes, des pictogrammes, des illustrations, des identités visuelles complètes. Ces compétences sont essentielles, mais elles ne suffisent pas : il faut aussi maîtriser la communication visuelle, la méthodologie de projet, la préparation à l'impression. C'est l'objet des parties suivantes.

Dans le chapitre suivant, nous aborderons la \terme{méthodologie de conception graphique}. Tu apprendras à comprendre le besoin du client, à rédiger et analyser un brief, à définir le public cible, à identifier le message principal, à déterminer les objectifs de communication, à rechercher des références, à construire un moodboard, à réaliser des esquisses, à développer plusieurs pistes créatives, à sélectionner et améliorer la proposition finale. Ces compétences te permettront de gérer un projet de design graphique de A à Z.

% ============================================================
\section*{Projet de la Partie III}
\addcontentsline{toc}{section}{Projet de la Partie III}
% ============================================================

\subsection*{Créer une identité visuelle complète}

\subsubsection*{Objectif}

Mettre en pratique l'ensemble des compétences acquises dans la Partie III : création de logos, illustration vectorielle, identité visuelle, typographie, couleurs, papeterie, présentation.

\subsubsection*{Sujet}

Créer une identité visuelle complète pour une marque fictive. Le sujet est libre, mais il doit être réaliste et cohérent. Exemples : café de quartier, marque de cosmétiques naturels, studio de yoga, agence de voyage, marque de vêtements éthiques, application mobile, etc.

\subsubsection*{Contraintes}

\begin{itemize}
    \item \terme{Logo} : création d'un logo professionnel (typographique, symbolique, mixte, abstrait).
    \item \terme{Variantes} : au moins 3 variantes (horizontale, verticale, symbole seul).
    \item \terme{Versions monochromes} : noir sur blanc, blanc sur noir.
    \item \terme{Palette} : 3 à 5 couleurs, avec codes (RVB, CMJN, Pantone).
    \item \terme{Typographies} : 1 à 2 familles, avec usages.
    \item \terme{Éléments secondaires} : icônes, pictogrammes, ou illustrations.
    \item \terme{Applications} : papeterie (carte de visite, papier à en-tête), web (bannière, avatar), packaging ou signalétique.
    \item \terme{Présentation} : document de présentation de l'identité visuelle.
\end{itemize}

\subsubsection*{Étapes}

\begin{enumerate}
    \item \terme{Brief} : définir la marque, le public, les valeurs, le positionnement.
    \item \terme{Recherche} : références, moodboard, esquisses.
    \item \terme{Création du logo} : esquisses, sélection, création dans Illustrator.
    \item \terme{Création de la palette} : couleurs, codes, usages.
    \item \terme{Choix des typographies} : familles, usages.
    \item \terme{Création des éléments secondaires} : icônes, pictogrammes, illustrations.
    \item \terme{Déclinaison} : papeterie, web, packaging ou signalétique.
    \item \terme{Présentation} : document de présentation, charte graphique.
\end{enumerate}

\subsubsection*{Auto-évaluation}

\begin{table}[h!]
\centering
\begin{tabular}{@{}lccccc@{}}
\toprule
\textbf{Critère} & \textbf{1} & \textbf{2} & \textbf{3} & \textbf{4} & \textbf{5} \\
\midrule
Brief compris et respecté & \(\square\) & \(\square\) & \(\square\) & \(\square\) & \(\square\) \\
Recherche pertinente & \(\square\) & \(\square\) & \(\square\) & \(\square\) & \(\square\) \\
Logo professionnel & \(\square\) & \(\square\) & \(\square\) & \(\square\) & \(\square\) \\
Variantes préparées & \(\square\) & \(\square\) & \(\square\) & \(\square\) & \(\square\) \\
Versions monochromes & \(\square\) & \(\square\) & \(\square\) & \(\square\) & \(\square\) \\
Palette cohérente & \(\square\) & \(\square\) & \(\square\) & \(\square\) & \(\square\) \\
Typographies adaptées & \(\square\) & \(\square\) & \(\square\) & \(\square\) & \(\square\) \\
Éléments secondaires & \(\square\) & \(\square\) & \(\square\) & \(\square\) & \(\square\) \\
Applications complètes & \(\square\) & \(\square\) & \(\square\) & \(\square\) & \(\square\) \\
Présentation claire & \(\square\) & \(\square\) & \(\square\) & \(\square\) & \(\square\) \\
Qualité globale & \(\square\) & \(\square\) & \(\square\) & \(\square\) & \(\square\) \\
\bottomrule
\end{tabular}
\caption{Grille d'auto-évaluation du projet de la Partie III.}
\end{table}

\subsubsection*{Livrables}

\begin{itemize}
    \item \terme{Fichier AI} : fichier source complet, avec calques organisés.
    \item \terme{Fichier PDF} : présentation de l'identité visuelle.
    \item \terme{Exports} : logo (SVG, PNG, JPEG), applications (PNG, PDF).
    \item \terme{Charte graphique} : document de règles d'utilisation.
\end{itemize}

\vspace{2em}

\begin{center}
\textit{Fin du chapitre 21 --- Illustration vectorielle et identité visuelle.}

\vspace{1em}

\textit{Fin de la Partie III --- Maîtrise d'Adobe Illustrator.}

\vspace{1em}

\textbf{Prochain chapitre : Chapitre 22 --- Méthodologie de conception graphique.}
\end{center}